%%
%% R-TypeRFC.tex for R-TypeRFC in /home/nerevarine/Documents/rtype-2015-tolsa_c/docs
%% 
%% Made by Nerevarine
%% Login   <nerevarine@epitech.net>
%% 
%% Started on  Sat Jan 19 22:20:20 2013 Nerevarine
%% Last update Sun Jan 20 16:56:41 2013 Nerevarine
%%

\documentclass[a4 paper, 12pt]{report}
\usepackage[english]{babel}
\usepackage[T1]{fontenc}
\usepackage[utf8x]{inputenc}
\usepackage{a4}
\usepackage[svgnames]{xcolor}
\usepackage{color}
\usepackage{fancybox}
\usepackage{fancyhdr}
\usepackage{graphicx}
\usepackage{tabularx}
\usepackage{multirow}
\usepackage{ucs}
\usepackage[colorlinks=true,urlcolor=SlateBlue]{hyperref}
\usepackage{url}
\usepackage{arydshln}

\newcommand{\email}[1]{\href{mailto:#1}{\nolinkurl{#1}}}
\newcommand{\actionsuccess}[1]{\textcolor{DarkGreen}{#1}}
\newcommand{\actionfailure}[1]{\textcolor{DarkRed}{#1}}


\title{Network Protocol For R-Type}
\author{Laurent Andral, Bertrand Boyer, Erick Clarivet\\
Nataniel Martin, Clément Petit, Camille Tolsa}
\date{\today}


%Entete et pieds de page.
\lhead{R-Type Group Project}
\chead{Network Protocol}
\rhead{\thepage}
\renewcommand{\headrulewidth}{1pt}
\renewcommand{\footrulewidth}{1pt}

\pagestyle{fancy}

\newcommand*{\plogo}{\includegraphics[height=2cm]{./Epitech.eps}} % Generic publisher logo

%-------------------------------------------------------
%TITLE PAGE
%----------------------------------------------------------

\newcommand*{\titleGM}{\begingroup % Create the command for including the title page in the document
\hbox{ % Horizontal box
\hspace*{0.2\textwidth} % Whitespace to the left of the title page
\rule{1pt}{\textheight} % Vertical line
\hspace*{0.05\textwidth} % Whitespace between the vertical line and title page text
\parbox[b]{0.75\textwidth}{ % Paragraph box which restricts text to less than the width of the page

{\noindent\Huge\bfseries Network Protocol \\[0.5\baselineskip] For R-Type}\\[2\baselineskip] % Title
{\large \textit{\today}}\\[4\baselineskip] % Tagline or further description
{\Large \textsc{Laurent Andral\\
Bertrand Boyer\\
Erick Clarivet\\
Nataniel Martin\\
Clément Petit\\
Camille Tolsa}} % Author name

\vspace{0.5\textheight} % Whitespace between the title block and the publisher
{\noindent \plogo}%\\[\baselineskip] % Publisher and logo
}}
\endgroup}

\begin{document}

%\maketitle

\thispagestyle{empty}
\titleGM

\tableofcontents
%\listoftables
\listoffigures

\chapter{Introduction}
This document describes the network protocol for R-Type Project.\\
It describes how the clients and the server will communicate together.\\
This protocol is only meant for the communation client-server type.

\subsection*{Status of this memo}
This document specifies a standart track protocol for the community of players. Distribution of this memo is unlimited.\\\\
\underline{Copyright Notice}: This project is Open-Source.

\subsection*{Conventions}
The key words ``MUST'', ``MUST NOT'', ``REQUIRED'', ``SHALL'', ``SHALL NOT'', ``SHOULD'', ``SHOULD NOT'', ``RECOMMENDED'', ``MAY'', and ``OPTIONAL'' in this document are to be interpreted as described in [RFC2119].

\subsection*{Internationalization considerations}
This document does not introduce or present any internationalization or localization issues.

\subsection*{Security considerations}
An injection of packets in UDP can harm the integrity of the system but it MUST provide the right MAGIC code and the right CHECKSUM hash.\\
And this is 0.0001\% likely to happen.


\chapter{Protocol packet}

\subsection*{Packet datagram}

\begin{figure}[!Hh]
  \centering
  \begin{tabular}{p{0.1cm}p{0.1cm}p{0.1cm}p{0.1cm}p{0.1cm}p{0.1cm}p{0.1cm}p{0.1cm}p{0.1cm}p{0.1cm}p{0.1cm}p{0.1cm}p{0.1cm}p{0.1cm}p{0.1cm}p{0.1cm}p{0.1cm}p{0.1cm}p{0.1cm}p{0.1cm}p{0.1cm}p{0.1cm}p{0.1cm}p{0.1cm}p{0.1cm}p{0.1cm}p{0.1cm}p{0.1cm}p{0.1cm}p{0.1cm}p{0.1cm}p{0.1cm}}
    \multicolumn{8}{c}{0} & \multicolumn{8}{c}{1} & \multicolumn{8}{c}{2} & \multicolumn{8}{c}{3}\\
    0 & 1 & 2 & 3 & 4 & 5 & 6 & 7 & 0 & 1 & 2 & 3 & 4 & 5 & 6 & 7 & 0 & 1 & 2 & 3 & 4 & 5 & 6 & 7 & 0 & 1 & 2 & 3 & 4 & 5 & 6 & 7\\
    \hline
    \multicolumn{8}{|c|}{0xFE} & \multicolumn{8}{c|}{0xE1} & \multicolumn{8}{c|}{0xDE} & \multicolumn{8}{c|}{0xAD}\\
    \hline
    \multicolumn{32}{|c|}{Checksum}\\
    \hline
    \multicolumn{16}{|c|}{Command ID} & \multicolumn{16}{c|}{Data size}\\
    \hline
    \multicolumn{32}{|c|}{Data}\\
    \multicolumn{32}{|c|}{}\\
    \hdashline
  \end{tabular}
  \caption{R-Type packet's datagram.}
\end{figure}

\subsection*{Description}

\begin{itemize}
\item The first 32 bits are used as a magic number to make sure the packet is conform (uint32\_t).
\item The next 32 bits, "Hashcode" to improve security of the packet and to check the packet integrity (uint32\_t).
\item The next 32 bits are used to define the timestamp of the packet, in order to check if the packet is to old, irrelevant (uint32\_t).
\item The next 16 bits are used for the ID command. It defines the incoming data pattern (uint16\_t).
\item The next 16 bits are used for the Data Size (uint16\_t).
\item The next XX bits defined by the data size are reserved for the strutured message.
\end{itemize}

\chapter{List of commands}
This chapter presents all commands authorized between client and server communication. All unknown command will be ignored.

\section{Client $\rightarrow$ server}
\subsection{Menu commands - TCP}
\begin{itemize}
\item AUTHENTIFICATION
  \begin{itemize}
  \item The client wants to authentify.
  \item ID of command: 000.
  \item Message: [Login:char[32]].
  \end{itemize}

\item SHOW\_ROOM
  \begin{itemize}
  \item Request of the list of rooms.
  \item ID of command: 001.
  \item Message: [empty].
  \end{itemize}

\item CREATE\_ROOM
  \begin{itemize}
  \item The client creates a new room.
  \item ID of command: 002.
  \item Message: [NameRoom:char[32]][NumberOfMaxPlayers:int][Difficulty:int][IDMap:int].
  \end{itemize}

\item JOIN\_ROOM
  \begin{itemize}
  \item The client joins the room.
  \item ID of command: 003.
  \item Message: [NameRoom:char[32]].
  \end{itemize}

\newpage

\item LEAVE\_ROOM
  \begin{itemize}
  \item The client leaves the room.
  \item ID of command: 004.
  \item Message: [empty].
  \end{itemize}

\item LAUNCH\_GAME
  \begin{itemize}
  \item The client, owner of the room, starts the game.
  \item ID of command: 005.
  \item Message: [NameRoom:char[32]].
  \end{itemize}
\end{itemize}

\subsection{In game commands - UDP}
\begin{itemize}
\item MOVEMENT
  \begin{itemize}
  \item The client moves.
  \item Id of command: 100.
  \item Message: [MovementType:int].
  \end{itemize}

\item SHOOTING
  \begin{itemize}
  \item The client shoots.
  \item ID of command: 101.
  \item Message: [IDGun:int].
  \end{itemize}

\item ASK\_DESCRIBE\_ENTITY
  \begin{itemize}
  \item The client ask to describe an entity.
  \item ID of command: 102.
  \item Message: [IDEntity:int].
  \end{itemize}

\item ASK\_PLAYER\_SCORE
  \begin{itemize}
  \item The client ask a players' score.
  \item ID of command: 103.
  \item Message: [IDEntity:int].
  \end{itemize}

\item ASK\_PLAYER\_LIFE
  \begin{itemize}
  \item The client ask a players' life.
  \item ID of command: 104.
  \item Message: [IDEntity:int].
  \end{itemize}
\end{itemize}

\section{Server $\rightarrow$ client}
\subsection{Menu commands - TCP}

\begin{itemize}
\item AUTHENTIFICATION\_OK
  \begin{itemize}
  \item Success of the authentification.
  \item ID of command: 200.
  \item Message: [empty].
  \end{itemize}

\item AUTHENTIFICATION\_KO
  \begin{itemize}
  \item Failure of the authentification.
  \item ID of command: 201.
  \item Message: [ErrorMessageID:int].
  \end{itemize}

\item ROOM\_DESCRIPTION
  \begin{itemize}
  \item The server sends a room description.
  \item ID of command: 202.
  \item Message: [NameRoom:char[32]][NumberOfMaxPlayers:int]
  \item[] [NumberOfCurrentPlayers:int][Difficulty:int][IDMap:int]
  \end{itemize}

\item ROOM\_PLAYER\_LIST
  \begin{itemize}
  \item The list of player's name in the room.
  \item ID of the command: 203.
  \item Message: [NameRoom:char[32]][Name1:char[32]]
  \item[] [Name2:char[32]][Name3:char[32]][Name4:char[32]]
  \end{itemize}

\item CREATE\_ROOM\_OK
  \begin{itemize}
  \item Success of the room creation.
  \item ID of command: 204.
  \item Message: [empty].
  \end{itemize}

\item CREATE\_ROOM\_KO
  \begin{itemize}
  \item Success of the room creation.
  \item ID of command: 205.
  \item Message: [ErrorMessageID:int].
  \end{itemize}

\newpage

\item JOIN\_ROOM\_OK
  \begin{itemize}
  \item Success of entering in the room.
  \item ID of command: 206.
  \item Message: [empty].
  \end{itemize}

\item JOIN\_ROOM\_KO
  \begin{itemize}
  \item Failure of entering in the room.
  \item ID of command: 207.
  \item Message: [ErrorMessageID:int].
  \end{itemize}

\item LEAVE\_ROOM\_OK
  \begin{itemize}
  \item Success of leaving the room.
  \item ID of command: 208.
  \item Message: [empty].
  \end{itemize}

\item LEAVE\_ROOM\_KO
  \begin{itemize}
  \item Failure of leaving the room.
  \item ID of command: 209.
  \item Message: [ErrorMessageID:int].
  \end{itemize}

\item LAUNCH\_GAME\_OK
  \begin{itemize}
  \item Success of launching the game.
  \item ID of command: 210.
  \item Message: [empty].
  \end{itemize}

\item LAUNCH\_GAME\_KO
  \begin{itemize}
  \item Failure of launching the game.
  \item ID of command: 211.
  \item Message: [ErrorMessageID:int].
  \end{itemize}
\end{itemize}

\subsection{In game commands - UDP}
\begin{itemize}
\item START\_GAME
  \begin{itemize}
  \item The server launch the server.
  \item ID of command: 300.
  \item Message: [empty].
  \end{itemize}

\item SPAWN\_ENTITY
  \begin{itemize}
  \item An entity has spawn on the map.
  \item ID of command: 301.
  \item Message: [IDEntity:int][IDType:int][PosX:int][PosY:int].
  \end{itemize}

\item DESTROY\_ENTITY
  \begin{itemize}
  \item An entity is destroyed.
  \item ID of command: 302.
  \item Message: [IDEntity:int].
  \end{itemize}

\item MOVE\_ENTITY
  \begin{itemize}
  \item An entity moves.
  \item ID of command: 303.
  \item Message: [IDEntity:int][PosX:int][PosY:int][MovementType:int].
  \end{itemize}

\item LIFE\_ENTITY
  \begin{itemize}
  \item The server send the life of the entity.
  \item ID of command: 304.
  \item Message: [IDEntity:int][Life:int].
  \end{itemize}

\item COLLISION
  \begin{itemize}
  \item Collision between two entities.
  \item ID of command: 305.
  \item Message: [IDEntity1:int][IDEntity2:int].
  \end{itemize}

\item DESCRIPTION\_ENTITY
  \begin{itemize}
  \item The server send the description of the entity.
  \item ID of command: 306.
  \item Message: [IDEntity:int][IDType:int][Life:int][PosX:int][PosY:int].
  \end{itemize}

\item ENTITY\_SCORE
  \begin{itemize}
  \item The server sends the score of the entity.
  \item ID of command: 307.
  \item Message: [IDEntity: int][Score:int].
  \end{itemize}

\item PLAYER\_DISCONNECT
  \begin{itemize}
  \item The server sends the player id who has just disconnected.
  \item ID of command: 308.
  \item Message: [IDEntity:int].
  \end{itemize}

\item END\_GAME
  \begin{itemize}
  \item The game has ended.
  \item ID of command: 309.
  \item Message: [MessageID:int].
  \end{itemize}
\end{itemize}


\chapter{Communication Client $\rightarrow$ Server}

This section presents actions and reactions from client and server for each packets sent. And it also gives the mains reasons for eache sending.

\section{TCP}

\subsection{Actions and reactions}

\begin{figure}[!Hh]
  \begin{itemize}
  \item \actionsuccess{Green} means the action succeed and \actionfailure{red} it failed.
  \item n is the number of room on the server (in game or not).
\end{itemize}

\centering
  \begin{tabular}{|l|l|}
    \hline
    \multicolumn{1}{|c|}{\textbf{Action}} & \multicolumn{1}{c|}{\textbf{Reaction}}\\
    \hline
    \hline
    \multirow{2}{*}{AUTHENTIFICATION} & \actionsuccess{AUTHENTIFICATION\_OK}\\
    & \actionfailure{AUTHENTIFICATION\_KO}\\
    \hline
    \multirow{2}{*}{SHOW\_ROOM} & n * ROOM\_DESCRIPTION\\
    & n * ROOM\_PLAYER\_LIST\\
    \hline
    \multirow{2}{*}{CREATE\_ROOM} & \actionsuccess{CREATE\_ROOM\_OK}\\
    & \actionfailure{CREATE\_ROOM\_KO}\\
    \hline
    \multirow{2}{*}{JOIN\_ROOM} & \actionsuccess{JOIN\_ROOM\_OK}\\
    & \actionfailure{JOIN\_ROOM\_KO}\\
    \hline
    \multirow{2}{*}{LEAVE\_ROOM} & \actionsuccess{LEAVE\_ROOM\_OK}\\
    & \actionfailure{LEAVE\_ROOM\_KO}\\
    \hline
    \multirow{2}{*}{LAUNCH\_GAME} & \actionsuccess{LAUNCH\_ROOM\_OK}\\
    & \actionfailure{LAUNCH\_ROOM\_KO}\\
    \hline
  \end{tabular}
  \caption{Actions and reactions in TCP.}
\end{figure}

\subsection{Reasons}

\begin{itemize}

\item AUTHENTIFICATION
  \begin{itemize}
  \item Is the login already exist ?
    \begin{itemize}
    \item Yes $\rightarrow$ send AUTHENFICATION\_KO.
    \item No $\rightarrow$ send AUTHENFICATION\_OK.
    \end{itemize}
  \end{itemize}

\item SHOW\_ROOM
  \begin{itemize}
  \item For each room, send ROOM\_DESCRIPTION to the client.
  \item For each room, send ROOM\_PLAYER\_LIST to the client.
  \end{itemize}

\item CREATE\_ROOM
  \begin{itemize}
  \item Is the room's name already exist ?
    \begin{itemize}
    \item Yes $\rightarrow$ send CREATE\_ROOM\_KO.
    \item No $\rightarrow$ send CREATE\_ROOM\_OK.
    \end{itemize}
  \end{itemize}

\item JOIN\_ROOM
  \begin{itemize}
  \item Is the room already in a game ?
    \begin{itemize}
    \item Yes $\rightarrow$ send JOIN\_ROOM\_KO.
    \item No. Is the room full ?
      \begin{itemize}
      \item Yes $\rightarrow$ send JOIN\_ROOM\_KO.
      \item No $\rightarrow$ send JOIN\_ROOM\_OK.
      \end{itemize}
    \end{itemize}
  \end{itemize}

\item LEAVE\_ROOM
  \begin{itemize}
    \item Send ROOM\_DESCRIPTION and ROOM\_PLAYER\_LIST to all players in the room.
    \item Is the client the last in the room ?
      \begin{itemize}
        \item Yes $\rightarrow$ Delete the room (server side).
      \end{itemize}
  \end{itemize}

\item LAUNCH\_ROOM
  \begin{itemize}
    \item Is the server ready for a new game ?
      \begin{itemize}
        \item Yes $\rightarrow$ send LAUNCH\_GAME\_OK to all players in the room.
          \item No $\rightarrow$ send LAUNCH\_GAME\_KO to the client.
      \end{itemize}
  \end{itemize}


\end{itemize}

\newpage

\section{UDP}

\subsection{Actions and reactions}

\begin{figure}[!Hh]
  Here is only presented actions which have an answer from the server.\\

  \centering
  \begin{tabular}{|l|l|}
    \hline
    \multicolumn{1}{|c|}{\textbf{Action}} & \multicolumn{1}{c|}{\textbf{Reaction}}\\
    \hline
    \hline
    \multirow{2}{*}{MOVEMENT} & \multirow{2}{*}{MOVE\_ENTITY}\\
    & \\
    \hline
    \multirow{2}{*}{SHOOTING} & \multirow{2}{*}{SPAWN\_ENTITY}\\
    & \\
    \hline
    \multirow{2}{*}{ASK\_DESCRIBE\_ENTITY} & \multirow{2}{*}{DESCRIPTION\_ENTITY}\\
    & \\
    \hline
    \multirow{2}{*}{ASK\_PLAYER\_SCORE} & \multirow{2}{*}{ENTITY\_SCORE}\\
    & \\
    \hline
    \multirow{2}{*}{ASK\_PLAYER\_PLAYER\_LIFE} & \multirow{2}{*}{LIFE\_ENTITY}\\
    & \\
    \hline
  \end{tabular}
  \caption{Actions and reactions in UDP.}
\end{figure}


\begin{figure}[!Hh]
  These events could be sent by the server without any request from the client.\\

  \centering
  \renewcommand{\arraystretch}{1.5}
  \begin{tabularx}{15cm}{|l|X|}
    \hline
    \multicolumn{1}{|c|}{\textbf{Event}} & \multicolumn{1}{c|}{\textbf{Reason}}\\
    \hline
    \hline
    SPAWN\_ENTITY & A new entity spawn on the map.\\
    \hline
    DESTROY\_ENTITY & An entity disapears or is killed.\\
    \hline
    MOVE\_ENTITY & An entity is moving (e.g. a player is pressing a movement key).\\
    \hline
    LIFE\_ENTITY & An ennemy or the player has been shot.\\
    \hline
    COLLISION & Two entity are touching each other.\\
    \hline
    DESCRIPTION\_ENTITY & After a spawned entity, the server can send this entity description.\\
    \hline
    ENTITY\_SCORE & After an ennemy is killed, the new score is sent to display.\\
    \hline
    PLAYER\_DISCONNECTED & The player leaves the game, the client doesn't need to display it anymore.\\
    \hline
  \end{tabularx}
  \caption{Events sent by the server.}
\end{figure}


\newpage

\section{Order of actions}
This part explains the order of actions for a game.

\begin{figure}[!Hh]
  \begin{center}
    \includegraphics[height=13cm]{./ClientInRoom.eps}
    \caption{Client side, before the game.}
  \end{center}
\end{figure}

\chapter{References}

[RFC 2119]
\newline
[RFC 2223]

\chapter{Authors adresses}

\begin{flushleft}
\begin{tabular}{lll}
  Laurent Andral & (andral\_l) & \email{laurent.andral@epitech.eu}\\
  Bertrand Boyer & (boyer\_b) & \email{bertrand.boyer@epitech.eu}\\
  Erick Clarivet & (clariv\_e) & \email{erick.clarivet@epitech.eu}\\
  Nataniel Martin & (mart\_n) & \email{nataniel.martin@epitech.eu}\\
  Clément Petit & (petit\_c) & \email{clement.petit@epitech.eu}\\
  Camille Tolsa & (tolsa\_c) & \email{camille.tolsa@epitech.eu}\\
\end{tabular}

\vspace{0.7cm}

R-Type Group Project\\
Epitech Toulouse\\
19 rue Bayard, 31000 Toulouse, France\\

\end{flushleft}


\end{document}
